\documentclass[10pt,letterpaper]{article}

\usepackage[T1]{fontenc}
\usepackage{lmodern}
\usepackage[margin=0.60in]{geometry}
\usepackage{array}
\usepackage{booktabs}
\usepackage{tabularx}
\usepackage{enumitem}
\usepackage[hidelinks]{hyperref}

\setlength{\parindent}{0pt}
\setlength{\parskip}{3pt}
\setlist[itemize]{leftmargin=1.25em,itemsep=1pt,topsep=2pt}
\renewcommand{\arraystretch}{1.12}
\newcommand{\openstatus}{\textbf{Open}}
\newcommand{\partialstatus}{\textbf{Partly supported}}

\begin{document}
\small

\begin{center}
{\Large\bfseries Remote Environmental Survey and Mapping of an Unseen Room}\\[3pt]
{\large Stage 1 Research Draft --- Work in Progress}\\[4pt]
Julian Joseph Amaro \quad Project I (247-519-VA) \quad Vanier College\\
October 23, 2026
\end{center}

\noindent\fbox{\parbox{0.96\linewidth}{
\textbf{Draft status.} This document records the research and project information available so far. It is not a claim that the rover has passed its requirements. No project-specific performance test has yet been recorded. The bill of materials and cost plan are intentionally excluded because they are separate Week 4 work.
}}

\section*{1. Change Log}

The submitted proposal remains the frozen baseline for needs N1--N8, watch constraints C1--C5, requirements REQ-RS-01--08, and risks R1--R3. The following additions make later project constraints explicit.

\begin{tabularx}{\linewidth}{>{\raggedright\arraybackslash}p{0.16\linewidth}X>{\raggedright\arraybackslash}p{0.16\linewidth}}
\toprule
\textbf{Item} & \textbf{Change and reason} & \textbf{Record} \\
\midrule
REQ-RS-09 & Added the exact Waveshare ESP32-S3-Touch-AMOLED-2.06 as a required output that displays the map, summary, and anomaly indication. & L-02, L-03.1 \\
REQ-RS-10 & Added the final custom-PCB requirement. Stage 1 may still use development modules, breadboard, perfboard, or protoboard. & L-02, L-03.1 \\
Gas interpretation & Clarified that the BME688 is a relative environmental/anomaly sensor, not selective toxic-gas detection or proof that a room is safe. & L-02--L-05, L-03.1 \\
\bottomrule
\end{tabularx}

A Raspberry Pi Zero 2 W and a 3D-printed tracked tank are candidate components only. Neither is a frozen architecture decision or a verified compatible choice.

\section*{2. Concept Alternatives}

Two genuinely different control approaches are under consideration.

\begin{tabularx}{\linewidth}{>{\bfseries}p{0.08\linewidth}X X}
\toprule
 & \textbf{Alternative A: centralized control} & \textbf{Alternative B: divided control} \\
\midrule
Approach & One main computer coordinates video, mapping, sensing, communications, and motor commands through a motor driver. & A main computer handles video, mapping, storage, and communications, while a microcontroller handles motors, encoders, and time-sensitive sensor input. \\
Main uncertainty & Whether one processor and operating system can provide sufficiently predictable motor, encoder, video, and mapping behaviour together. & Whether the extra controller, firmware, interfaces, and synchronization are justified by improved real-time control. \\
Current status & Not scored; Raspberry Pi Zero 2 W remains a candidate. & Not scored; the microcontroller and interface are not selected. \\
\bottomrule
\end{tabularx}

The evaluation criteria and weights must be chosen and dated in the logbook before either alternative is scored. Likely criteria include requirement coverage, mapping risk, real-time control, implementation effort, power demand, integration complexity, and evidence obtainable during Project I. No winner is claimed yet.

\textbf{Functional flow under investigation:}

\begin{center}
\fbox{Remote command / live view}
$\longleftrightarrow$
\fbox{Motion and position estimation}
$\longrightarrow$
\fbox{Environmental sensing}\\[4pt]
$\longrightarrow$
\fbox{Map and tagged-record generation}
$\longrightarrow$
\fbox{Compact summary transfer}
$\longrightarrow$
\fbox{Watch display}
\end{center}

\section*{3. Feasibility Case}

Claims begin at C6 because C1--C5 already identify watch constraints in the frozen proposal.

\begin{tabularx}{\linewidth}{>{\bfseries}p{0.055\linewidth}>{\raggedright\arraybackslash}p{0.12\linewidth}X>{\raggedright\arraybackslash}p{0.18\linewidth}>{\raggedright\arraybackslash}p{0.12\linewidth}}
\toprule
ID & Tests & Claim and present evidence & Class / source & Status \\
\midrule
C6 & RS-06, RS-09 & The required watch has a 410 by 502 AMOLED touch display, ESP32-S3R8, 8 MB PSRAM, 32 MB flash, Wi-Fi b/g/n, and BLE 5. These documented capabilities make a compact map and survey-summary interface plausible, but the project UI and 5 s transfer are untested. & Documented specifications; L-02 \cite{waveshare,espressif} & \partialstatus \\
C7 & RS-03, RS-04 & The BME688 documents temperature, humidity, pressure, and broad gas response. It can support a relative anomaly channel, but it is not selective or certified toxic-gas detection. & Documented research; L-02, L-03 \cite{bosch} & \partialstatus \\
C8 & RS-04 & Public BME688 projects demonstrate logging and experimental classification workflows, but no inspected public data set covered methane, propane/butane, CO, and H$_2$S together. No published result was reproduced. & Literature/repository review; L-05 & \openstatus \\
C9 & RS-05 & The tracked tank may provide a mechanical base, but its encoder integration, payload capacity, track slip, and mapping accuracy are not established for this project. & Documented/inferred; L-02 & \openstatus \\
C10 & RS-01, RS-02, RS-05 & A Pi Zero 2 W may be able to handle video and mapping, but processor load, thermal behaviour, motor-control timing, and power demand have not been calculated or measured for the final architecture. & Assumed; planned calculation and test & \openstatus \\
C11 & RS-01--RS-10 & The integrated rover can meet the complete requirement set. No integrated hardware test exists yet. & Assumed; tests pending & \openstatus \\
\bottomrule
\end{tabularx}

The research supports technical plausibility, not requirement compliance. In particular, no 2 m drive, video-latency test, mapping trial, BME688 bench test, watch-transfer timing, power test, or clean build has been completed. Vendor specifications are documented evidence about components, not measured project results.

\section*{4. Baselined Requirements}

\begin{tabularx}{\linewidth}{>{\bfseries}p{0.12\linewidth}X>{\raggedright\arraybackslash}p{0.18\linewidth}}
\toprule
\textbf{ID} & \textbf{Current requirement and reason} & \textbf{Status} \\
\midrule
REQ-RS-01 & Complete a 2 m remotely commanded path with the operator outside the room; demonstrates remote entry. & Open: drive test pending. \\
REQ-RS-02 & Display live phone video with latency at most 1 s; supports obstacle response. & Open: five timing trials pending. \\
REQ-RS-03 & Map at least three tagged temperature, humidity, and relative-gas records; demonstrates spatial variation. & Partly supported: sensor capability documented; integration open. \\
REQ-RS-04 & Flag a potential anomaly at 125\% or more of baseline; this is a relative indicator, not a safety limit. & Partly supported: rule defined; implementation and sensor behaviour open. \\
REQ-RS-05 & Map walls within $\pm$10 cm in at least four of five runs; preserves the proposal accuracy target. & Open: highest-priority mapping test pending. \\
REQ-RS-06 & Deliver the completed summary to the watch within 5 s; provides portable results. & Partly supported: radio capability documented; transfer time open. \\
REQ-RS-07 & Keep total hardware cost at or below CAD \$700; stakeholder constraint N7. & Open: cost work is separate and excluded here. \\
REQ-RS-08 & Build firmware from a clean copy using one documented command; supports reproducibility. & Open: clean-build demonstration pending. \\
REQ-RS-09 & Use the exact Waveshare watch to display the map, summary, and anomaly indications. & Partly supported: hardware is suitable in principle; UI unbuilt. \\
REQ-RS-10 & Include a custom PCB in the final system; breadboard/modules alone are insufficient for the final build. & Open: binding plan; PCB not yet designed or built. \\
\bottomrule
\end{tabularx}

\section*{5. Risk Register}

The following are the present research-backed risks. Probability and consequence scores are not assigned yet because no measured project evidence supports a defensible ranking. The required $P=p\times c$ matrix remains open.

\begin{tabularx}{\linewidth}{>{\bfseries}p{0.06\linewidth}X X}
\toprule
ID & \textbf{Risk, assumption, and consequence} & \textbf{Evidence needed and mitigation} \\
\midrule
R1 & Wheel slip and range noise may exceed the 10 cm error limit. Assumes odometry and ranging remain repeatable on the selected chassis. Failure would make wall and reading locations unreliable. & Five mapping runs against eight tape-measured points. Calibrate geometry, reduce speed, add encoders/ranging, or change mapping method. \\
R2 & BME688 response may drift with humidity or be too slow for brief stops. Assumes a stable relative baseline. Failure could produce misleading anomaly indications. & Controlled safe baseline, warm-up, humidity, dwell-time, and repeatability tests. Keep claims relative and change the sensor or sampling method if needed. \\
R3 & Doorway wireless traffic may delay phone video or the watch summary. Assumes the available 2.4 GHz links coexist under load. Failure would affect REQ-RS-02 and REQ-RS-06. & Measure video and summary delivery together. Reduce data rate, buffer/resend, or test a phone-mediated path. \\
R4 & The selected controller architecture may not meet processing, timing, thermal, or power needs. & Calculate the power/load budget and prototype the two highest-load functions before selection. Use divided control if centralized timing is inadequate. \\
\bottomrule
\end{tabularx}

\section*{6. Gate 1 Readiness}

The current research supports a \textbf{preliminary proceed-with-conditions} position. Official documentation supports the watch as a compact wireless display, and the BME688 can support environmental readings and a relative anomaly channel. Research also identifies important limits: the BME688 cannot certify safety or selectively identify household toxic gases, the tracked tank is not yet proven as a mapping platform, and the Pi Zero 2 W architecture has not been selected or validated.

The conditions are to fix and log the concept-selection criteria and weights, score the two architectures, select the controller structure, bench-test the BME688, measure mapping accuracy, video latency and watch delivery, verify a clean build, and calculate the power/current budget. Until those steps are complete, every performance requirement remains open or partly supported.

\begin{thebibliography}{9}
\small
\bibitem{waveshare}
Waveshare, ``ESP32-S3-Touch-AMOLED-2.06 Product Page and Wiki,'' accessed October 23, 2026.

\bibitem{espressif}
Espressif Systems, ``ESP32-S3 Series Datasheet'' and ``RF Coexistence Guide,'' accessed October 23, 2026.

\bibitem{bosch}
Bosch Sensortec, ``BME688 Gas Sensor Product Information,'' accessed October 23, 2026.
\end{thebibliography}

\end{document}
